\documentclass{amsart}

% --- PACKAGES ---
% Provides control over page margins
\usepackage[margin=1in]{geometry}
% Provides access to various mathematical symbols (loaded by amsart)
\usepackage{amssymb} 
% Provides clickable internal and external links
\usepackage[colorlinks=true, allcolors=black, pdftex]{hyperref}

% --- DOCUMENT METADATA (for PDF properties) ---
\hypersetup{
    pdftitle={Polylogarithm Values at a Golden Ratio-Based Argument},
    pdfauthor={Tristen Harr},
    pdfsubject={Complex Analysis, Special Functions, Number Theory},
    pdfkeywords={Infinite Series, Polylogarithm, Golden Ratio, Quasicrystals, Special Functions, Numerical Analysis, Transcendental Number Theory}
}

% --- THEOREM-LIKE ENVIRONMENTS ---
% 'plain' style for theorems, lemmas, etc. (italicized text)
\theoremstyle{plain}
\newtheorem{theorem}{Theorem}[section]
\newtheorem{proposition}[theorem]{Proposition}
\newtheorem{lemma}[theorem]{Lemma}
\newtheorem{corollary}[theorem]{Corollary}
\newtheorem{conjecture}[theorem]{Conjecture}

% 'definition' style (upright text)
\theoremstyle{definition}
\newtheorem{definition}[theorem]{Definition}

% 'remark' style (upright text, no numbering by default)
\theoremstyle{remark}
\newtheorem*{remark}{Remark}

% --- CUSTOM OPERATORS AND COMMANDS ---
% For the Polylogarithm function operator
\DeclareMathOperator{\Li}{Li} 
% Custom command for the main constant for consistency and ease of change
\newcommand{\LambdaG}{\Lambda_{G1}}
% Custom command for the golden ratio symbol
\newcommand{\phiG}{\phi}

% ======================================================================
% --- BEGIN DOCUMENT ---
% ======================================================================

\begin{document}

\title[Polylogarithms at a Golden Ratio Argument]{Polylogarithm Values at a Golden Ratio-Based Argument}
\author{Tristen Harr}
\date{Revised June 2024}

\begin{abstract}
This paper introduces and analyzes a specific complex constant, denoted $\LambdaG$, which is derived from inverse powers of the golden ratio, $\phiG$. We define the constant's real and imaginary components and prove that its magnitude is less than one. This validates the use of $\LambdaG$ as an argument in the Polylogarithm function, $\Li_s(z)$. We present high-precision numerical evaluations for the Dilogarithm ($s=2$) and Trilogarithm ($s=3$) of this constant. Based on this numerical evidence, we conjecture that the values $\Li_s(\LambdaG)$ are transcendental for all integers $s \ge 2$. Furthermore, we conjecture that these values do not belong to the field extension $\mathbb{Q}(\pi, \ln(2), \phiG)$, a field chosen based on known polylogarithm identities. This line of inquiry is motivated in part by potential applications in the study of quasicrystals.
\end{abstract}

\maketitle

% --- SECTION 1: INTRODUCTION ---
\section{Introduction}
The study of special functions and constants often reveals connections between different areas of mathematics. This paper focuses on the analysis of the Polylogarithm function at a novel complex argument, $\LambdaG$, derived from the golden ratio, $\phiG = (1+\sqrt{5})/2$. A review of foundational literature, such as Lewin's \textit{Polylogarithms and Associated Functions}~\cite{Lewin}, reveals no prior study of this specific complex argument. To substantiate this claim, a search of major academic databases (including MathSciNet and the arXiv preprint server, conducted up to June 2024) was performed, which likewise showed no indication of this particular argument being studied, reinforcing the novelty of this investigation.

The motivation for this constant stems from the classical problem of dividing a line segment in extreme and mean ratio, as found in Book VI, Proposition 30 of Euclid's \textit{Elements}~\cite{Euclid}. The paper proceeds by rigorously defining $\LambdaG$ from inverse powers of $\phiG$ and proving its validity as a polylogarithm argument. The analysis of the resulting values, $\Li_s(\LambdaG)$, forms the central topic of this work. Understanding the nature of these specific values may offer insights into mathematical domains where the golden ratio is foundational. For instance, the diffraction patterns of quasicrystals are intrinsically linked to the golden ratio~\cite{Levine1984}, and polylogarithms frequently arise in quantum field theory calculations. Because physical properties of such systems (e.g., electron transport in aperiodic structures, where models often involve complex hopping parameters) involve complex phase factors, it is plausible that a constant like $\LambdaG$—which combines the golden ratio with a non-trivial complex phase in a geometrically-grounded way—could serve as a more natural argument for these functions than purely real values. This work explores the possibility that $\Li_s(\LambdaG)$ might emerge as a fundamental constant in the theoretical description of aperiodic structures. The paper culminates in a formal conjecture about the transcendental nature of these values, based on high-precision numerical analysis.

% --- SECTION 2: THE CONSTANT AND ITS PROPERTIES ---
\section{Definition, Properties, and Polylogarithm Convergence}

\subsection{Defining Components from the Golden Ratio}
We define two real constants, $T$ and $J$, directly from the algebraic properties of the golden ratio.
\begin{definition}[Primary Real Components]
The constants $T$ and $J$ are defined from the golden ratio, $\phiG$, as:
\[ T = \frac{1}{2\phiG} \quad \text{and} \quad J = \frac{1}{2\phiG^2} \]
\end{definition}

These definitions lead to a pair of linear relations that highlight their connection to classical geometry; as will be further elaborated in Appendix A, they are chosen to correspond directly to the division of a line segment in the golden ratio.
\begin{proposition}[System of Equations]
The constants $T$ and $J$ are the unique real numbers $(a,b)$ that satisfy the system:
\begin{align}
    a/b &= \phiG \label{eq:1} \\
    a+b &= 1/2 \label{eq:2}
\end{align}
\end{proposition}
\begin{proof}
First, we show that $T$ and $J$ satisfy the system. The ratio is immediate from the definitions: $T/J = (1/(2\phiG))/(1/(2\phiG^2)) = \phiG^2/\phiG = \phiG$. For the sum, we use the identity $\phiG+1 = \phiG^2$:
\[ T+J = \frac{1}{2\phiG} + \frac{1}{2\phiG^2} = \frac{\phiG+1}{2\phiG^2} = \frac{\phiG^2}{2\phiG^2} = \frac{1}{2} \] 
The two equations, (\ref{eq:1}) and (\ref{eq:2}), form a linear system in variables $(a,b)$, which can be written as $a - \phiG b = 0$ and $a+b=1/2$. The determinant of the coefficient matrix is $\begin{vmatrix} 1 & -\phiG \\ 1 & 1 \end{vmatrix} = 1 - (-\phiG) = 1+\phiG$. Since $1+\phiG = \phiG^2 \neq 0$, the system has a unique solution.
\end{proof}

\begin{remark}
This formulation shows that the constraint $T+J=1/2$ is not an arbitrary choice, but a direct consequence of defining the components from successive inverse powers of $\phiG$. Using $\phiG = (1+\sqrt{5})/2$, the explicit values are $T = (\sqrt{5}-1)/4$ and $J = (3-\sqrt{5})/4$.
\end{remark}

From these constants, we define our primary object of study.
\begin{definition}[The Golden Ratio-Based Constant \texorpdfstring{$\LambdaG$}{Lambda\_G1}]
The constant $\LambdaG$ is defined as the complex number $\LambdaG = T+iJ$.
\end{definition}

\begin{remark}
The subscript '1' in $\LambdaG$ is used to distinguish this constant, based on the first and second inverse powers of $\phiG$, from a potential family of related constants, $\Lambda_{Gn}$, which could be constructed from higher inverse powers (e.g., $\phiG^{-n}$ and $\phiG^{-(n+1)}$). The investigation of such a family is a topic for future work.
\end{remark}

As an algebraic number, $\LambdaG$ is a root of an integer polynomial.
\begin{proposition}[Minimal Polynomial of \texorpdfstring{$\LambdaG$}{Lambda\_G1}]
The constant $\LambdaG$ is an algebraic number whose minimal polynomial over $\mathbb{Q}$ is $16x^4 + 16x^3 + 24x^2 - 20x + 5 = 0$.
\end{proposition}
\begin{proof}
Let $z = \LambdaG$. Using the explicit values $T = (\sqrt{5}-1)/4$ and $J = (3-\sqrt{5})/4$, we have $z = \frac{\sqrt{5}-1}{4} + i\frac{3-\sqrt{5}}{4}$. To eliminate the radicals and the imaginary unit, we sequentially isolate terms and square both sides of the equation.
\begin{align*}
    4z &= (\sqrt{5}-1) + i(3-\sqrt{5}) \\
    4z + 1 - \sqrt{5} &= i(3-\sqrt{5})
\end{align*}
Squaring both sides gives:
\begin{align*}
    (4z+1 - \sqrt{5})^2 &= -(3-\sqrt{5})^2 \\
    (4z+1)^2 - 2(4z+1)\sqrt{5} + 5 &= -(9 - 6\sqrt{5} + 5) \\
    16z^2 + 8z + 1 - 2(4z+1)\sqrt{5} + 5 &= -14 + 6\sqrt{5} \\
    16z^2 + 8z + 6 - (8z+2)\sqrt{5} &= -14 + 6\sqrt{5}
\end{align*}
Rearranging to isolate the $\sqrt{5}$ term:
\begin{align*}
    16z^2 + 8z + 20 &= (8z+2+6)\sqrt{5} \\
    16z^2 + 8z + 20 &= (8z+8)\sqrt{5} \\
    4(4z^2 + 2z + 5) &= 8(z+1)\sqrt{5} \\
    (4z^2 + 2z + 5) &= 2(z+1)\sqrt{5}
\end{align*}
Squaring both sides again eliminates the final radical:
\begin{align*}
    (4z^2 + 2z + 5)^2 &= 4(z+1)^2 \cdot 5 \\
    16z^4 + 16z^3 + 44z^2 + 20z + 25 &= 20(z^2 + 2z + 1) \\
    16z^4 + 16z^3 + 44z^2 + 20z + 25 &= 20z^2 + 40z + 20
\end{align*}
Combining terms yields the polynomial with integer coefficients:
\[ 16z^4 + 16z^3 + 24z^2 - 20z + 5 = 0 \]
This polynomial is irreducible over $\mathbb{Q}$, a fact verifiable with standard computational algebra systems (e.g., using the \texttt{is\_irreducible()} method in SageMath), and thus is the minimal polynomial of $\LambdaG$.
\end{proof}

\subsection{Convergence Property}
A crucial property of $\LambdaG$ is that its magnitude is less than one, which is required for the convergence of the Polylogarithm series.
\begin{lemma}[Magnitude of $\LambdaG$]
\label{lem:convergence}
The constant $\LambdaG$ lies within the unit disk, i.e., $|\LambdaG| < 1$.
\end{lemma}
\begin{proof}
We require $|\LambdaG|^2 < 1$. We calculate the squared values of the components, $T^2 = \frac{6-2\sqrt{5}}{16}$ and $J^2 = \frac{14-6\sqrt{5}}{16}$. Summing these gives:
\begin{align*}
    |\LambdaG|^2 &= T^2+J^2 = \frac{6-2\sqrt{5} + 14-6\sqrt{5}}{16} \\
                 &= \frac{20-8\sqrt{5}}{16} = \frac{5-2\sqrt{5}}{4} \approx 0.131966
\end{align*}
Since $0.131966 < 1$, the condition is satisfied.
\end{proof}

\subsection{The Polylogarithm at \texorpdfstring{$\LambdaG$}{Lambda\_G1}}
The property $|\LambdaG| < 1$ is central to this paper, as it allows us to define the Polylogarithm values at this constant.
\begin{definition}[The Polylogarithm Function]
For any complex $s$, the Polylogarithm function of order $s$, $\Li_s(z)$, is defined by the power series:
\[ \Li_s(z) = \sum_{n=1}^{\infty} \frac{z^n}{n^s} \]
This series converges for all complex numbers $z$ such that $|z|<1$.
\end{definition}

Since $\LambdaG$ satisfies the convergence condition, we can define the specific numerical values that are the subject of our main conjecture:
\[ \Li_s(\LambdaG) = \sum_{n=1}^{\infty} \frac{(\LambdaG)^{n}}{n^s} \]
This identification holds for any complex $s$.

% --- SECTION 3: NUMERICAL ANALYSIS AND A CONJECTURE ---
\section{Numerical Analysis and a Conjecture on the Nature of \texorpdfstring{$\Li_s(\LambdaG)$}{Li\_s(Lambda\_G1)}}

\subsection{High-Precision Numerical Evaluations}
We present high-precision numerical evaluations for the first two non-trivial integer cases of the Polylogarithm of $\LambdaG$. The values were computed using the Python \texttt{mpmath} library~\cite{mpmath} with a working precision of 200 decimal digits to ensure robustness.
\begin{itemize}
    \item For $s=2$, the Dilogarithm of $\LambdaG$ is:
    \[ \Li_2(\LambdaG) \approx 0.322328253720993 + 0.226730769307749\;i \]
    
    \item For $s=3$, the Trilogarithm of $\LambdaG$ is:
    \[ \Li_3(\LambdaG) \approx 0.310345112613931 + 0.198884393433621\;i \]
\end{itemize}

\subsection{A Conjecture on the Nature of These Values}
The transcendental nature of special function values at algebraic arguments is a central topic in number theory~\cite{Baker}. For context, values of the Dilogarithm are known to be transcendental at certain related real algebraic points, such as $\Li_2(1/2) = \frac{\pi^2}{12} - \frac{1}{2}\ln^2(2)$ and $\Li_2(\phiG^{-2}) = \frac{\pi^2}{15} - \ln^2(\phiG)$. The constant $\LambdaG$, being a non-real algebraic number that is not a root of unity, presents a distinct and challenging case. Transcendence proofs, including those based on Baker's theory, frequently rely on techniques that do not readily extend to a general complex algebraic argument like $\LambdaG$. The constant $\LambdaG$ falls into none of the categories (rational, real, root of unity) for which the most powerful results are available, placing its analysis outside the scope of many established methods.

To investigate the nature of these values, we employed the PSLQ integer relation algorithm~\cite{Bailey2007}. The real and imaginary components of $\Li_2(\LambdaG)$ and $\Li_3(\LambdaG)$ were each tested for the existence of an integer relation against a basis of well-known mathematical constants. This basis was chosen specifically because the constants $\pi$, $\ln(2)$, and $\ln(\phiG)$ (which is related to $\phiG$ over $\mathbb{Q}$) are the precise components of known closed-form evaluations such as $\Li_2(\phiG^{-2}) = \frac{\pi^2}{15} - \ln^2(\phiG)$, making this the most plausible field to test for relations. The basis was composed of elements of the form $\pi^a (\ln(2))^b \phiG^c$ where $a,b,c$ are non-negative integers with total degree $a+b+c < 5$. The computational cost of PSLQ searches increases significantly with the size of the basis; the chosen degree limit represents a practical balance between search breadth and computational feasibility for achieving the stated precision. No relation was found with a detection tolerance of $10^{-190}$. While not a formal proof, this failure provides strong numerical motivation for their apparent transcendental nature and leads us to propose the following general conjecture.

\begin{conjecture}
For any integer $s \ge 2$, the value $\Li_s(\LambdaG)$ is a transcendental number. Furthermore, it is conjectured that this number is not in the field $\mathbb{Q}(\pi, \ln(2), \phiG)$.
\end{conjecture}

\begin{remark}
Proving this conjecture is beyond the scope of this paper, but proposing it frames a clear direction for future work. The field in the second part of the conjecture is a well-motivated but non-exhaustive choice, intended to precisely formulate the negative result of our numerical search.
\end{remark}

% --- SECTION 4: CONCLUSION ---
\section{Conclusion}
This paper has introduced a complex constant, $\LambdaG$, whose definition is derived non-arbitrarily from the golden ratio. We have proven that it is an algebraic number lying within the unit disk, which validates its use as an argument in the Polylogarithm function. The primary contribution of this work is the introduction of this constant and the subsequent analysis of its polylogarithmic values, $\Li_s(\LambdaG)$. Based on high-precision numerical evidence from the $s=2$ and $s=3$ cases~\cite{mpmath}, we have proposed a new open problem: a conjecture on the transcendental nature of these values.

% --- SECTION 5: FUTURE WORK ---
\section{Future Work}
The conjecture presented in Section 3 opens a clear path for future investigation. A primary goal remains to prove or disprove the transcendental nature of $\Li_s(\LambdaG)$. More specific avenues for future work include:
\begin{itemize}
    \item Investigating further consequences of the algebraic properties of $\LambdaG$. Its Galois conjugates consist of its complex conjugate and another complex conjugate pair, $\frac{-1\pm\sqrt{5}}{4} \pm i\frac{3\mp\sqrt{5}}{4}$. Exploring whether relations exist between the polylogarithm values at these conjugate points could be fruitful.
    \item Separately analyzing the real and imaginary parts of $\Li_s(\LambdaG)$ and determining if they are transcendental or satisfy relations with other known constants.
    \item Expanding the numerical search for relations by using a larger basis of constants (e.g., including Euler's constant $\gamma$ and Catalan's constant) to test for more complex identities, alongside an increased search depth.
    \item Investigating potential applications in physics, particularly in the study of quasicrystals. Given the link between the golden ratio and diffraction patterns, and the appearance of polylogarithms in quantum field theory calculations, these values may be fundamental in describing aperiodic structures.
    \item Exploring related series identities, such as the evaluation using Eulerian polynomials shown in Appendix B, to see if they reveal further a structure or properties of $\Li_s(\LambdaG)$.
\end{itemize}

% --- APPENDICES AND REFERENCES ---
\appendix
\section{Geometric Origin of the Constants \texorpdfstring{$T$ and $J$}{T and J}}
The definitions of $T$ and $J$ have a direct geometric motivation arising from the division of a line segment according to the golden ratio. If a segment of length $L$ is divided into two pieces, $a$ and $b$, such that $a/b = \phiG$, the pieces have lengths $a=L/\phiG$ and $b=L/\phiG^2$. Our constants $T$ and $J$ correspond to this construction for $L=1/2$, providing a non-arbitrary basis for the constant's definition. The reader may find it helpful to visualize this construction: a line segment of length $1/2$ is composed of a longer part $T$ and a shorter part $J$.

\section{Evaluation of a Related Polynomial-Weighted Series}
As a consequence of Lemma~\ref{lem:convergence}, $\LambdaG$ is a valid argument in the series whose sum is given by Eulerian polynomials, $A_k(x)$. This provides a direct evaluation for a family of series with $\LambdaG$ as the base. For any integer $k \ge 0$:
\[ \sum_{n=0}^{\infty} n^k \cdot (\LambdaG)^n = \frac{A_k(\LambdaG)}{(1-\LambdaG)^{k+1}} \]
While not central to the paper's main conjecture, this identity is an interesting property of the constant and demonstrates its utility in evaluating other classes of well-known series.

\begin{thebibliography}{9}

\bibitem{Bailey2007}
D. H. Bailey. "The PSLQ Integer Relation Algorithm". In: \textit{Computational Science and Its Applications -- ICCSA 2007}. Springer, 2007, pp. 1--14.

\bibitem{Baker}
A. Baker. \textit{Transcendental Number Theory}. Cambridge University Press, 1975. 

\bibitem{Euclid}
T. L. Heath, translator. \textit{The Thirteen Books of Euclid's Elements}, 2nd ed. Dover Publications, 1956. 

\bibitem{mpmath}
F. Johansson, et al. \textit{mpmath: a Python library for arbitrary-precision floating-point arithmetic}, Version 1.3.0, \url{http://mpmath.org/}, 2023. (Accessed June 18, 2024). 

\bibitem{Levine1984}
D. Levine and P. J. Steinhardt. "Quasicrystals: A New Class of Ordered Structures". \textit{Physical Review Letters}, 53(26), 2477--2480, 1984. 

\bibitem{Lewin}
L. Lewin. \textit{Polylogarithms and Associated Functions}. North-Holland, 1981.

\end{thebibliography}

\end{document}